%%%%%%%%%%  scan idx 39  |  folio 26  %%%%%%%%%%
premier $\lambda$-anneau défini ici, si la caractéristique de $k$ est $0$, mais non dans le cas général.

3) \emph{Détermination de $\underline{K}(G)$ et de $\underline{K}_r(G)$ dans certains cas importants}~: On suppose le corps $k$ algébriquement clos, le groupe $G$ algébrique \emph{affine connexe} et «~globalement \emph{réductif}~» (le radical de $G$ est isomorphe à $k^{*s}$) et l'on ne considère que les représentations linéaires rationnelles. En caractéristique $0$, la complète réductibilité des représentations linéaires rationnelles de $G$ montre que $\underline{K}(G)=\underline{K}_r(G)$ est le groupe libre engendré par les classes des réprésentations rationnelles irréductibles. Or, si $T$ est un tore maximal de $G$, on a $\underline{K}(T)=\underline{K}_r(T)=\Z(\That)$, algébre du groupe $\That$ des caractères rationnels de $T$~; on a $\That\simeq\Z^r$ si $r$ est le rang de $G$ (ceci est indépendant de la caractéristique). Le groupe de Weyl $W$ opère dans $T$, donc dans $\underline{K}(T)$, d'une manière qui ne dépend que des opérations de $W$ dans le groupe libre $\That$. On a un homomorphisme évident de $\lambda$-anneaux, déduit de l'injection de $T$ dans $G$~:
\begin{equation}\tag{1.13}
\underline{K}_r(G)\longrightarrow \underline{K}_r(T)^W=\underline{K}(T)^W,
\end{equation}
où $\underline{K}(T)^W$ dénote l'ensemble des points fixes de $W$ dans $\underline{K}(T)$. La théorie des représentations linéaires de $G$\footnote{Cf.\ Séminaire Chevalley 1956/58 Groupes de Lie algébriques (École Normale Supérieure).} démontre alors le théorème suivant~:

\textbf{Théorème 1.1.} \emph{L'homomorphisme} (1.13) \emph{est un isomorphisme.}

On voit donc que $\underline{K}_r(G)$, une fois connu le «~type~» de $G$ (caractérisé par $W$ opérant sur un groupe $\Z^r$), ne dépend pas de la caractéristique, et que c'est un $\lambda$-anneau spécial. On a une base canonique de $\underline{K}_r(G)$ correspondant aux orbites de $W$ dans $\That$~; une autre base s'obtient ainsi~:
